\documentclass[10pt,letterpaper]{article}

% =========================================================
% PAGE AND FONT SETUP
% =========================================================
\usepackage[
	letterpaper,
	left=1in,
	right=1in,
	top=1in
]{geometry}

\usepackage[T1]{fontenc}
\usepackage[utf8]{inputenc}
\usepackage{newtxtext}
\usepackage[hidelinks]{hyperref}
\usepackage{microtype}
\usepackage{xurl}

\pagestyle{empty}
\setlength{\parindent}{0pt}
\setlength{\parskip}{0.85em}
\urlstyle{same}

\begin{document}

\begin{flushleft}
Nahid Hasan Sagar\\
Bloomington, Indiana\\
+1 (930) 904-1643\\
nahidhasansagar@gmail.com
\end{flushleft}

July 31, 2026

Hiring Committee\\
Department of Pathology and Laboratory Medicine\\
Indiana University and IU Health Physicians\\
Indianapolis, Indiana

Dear Hiring Committee,

I am writing to apply for the Research Data Analyst position in the Department of Pathology and Laboratory Medicine. As an M.S. candidate in Statistical Science at Indiana University Bloomington, I have built a strong foundation in computational statistics, reproducible data analysis, and benchmarking research workflows. The opportunity to support a computational team working with machine learning, image analysis, and research data pipelines strongly matches my training and interests.

In my current role as a Research Assistant in the StatMIND Lab, I develop efficient C++ code with Rcpp and RcppArmadillo for high-dimensional neuroimaging workflows, reproduce original results from R implementations, and validate correctness through unit testing and benchmarking. That work has required careful attention to data quality, consistency across runs, and clear presentation of results to technical audiences. I also built reproducible benchmarking workflows to compare runtime and scalability across implementations, and helped achieve a 4x to 5x speedup in core computations.

This experience is directly relevant to the responsibilities in your posting, including managing research data, benchmarking analysis pipelines, and presenting results to collaborators with different technical backgrounds. My coursework and thesis work in Bayesian modeling, time-series analysis, and statistical computing have further strengthened my ability to evaluate quantitative methods, write clear technical reports, and support data-driven research. In addition, my experience as a Teaching Assistant and long-term instructor has strengthened my communication, organization, and mentoring skills in fast-paced settings.

Along with my statistical training, I bring strong problem-solving skills from my physics background and childhood facination with computers. I can quickly learn new programs as well as any specialized software used in your department. I am confident that my combination of statistical expertise, programming experience, and commitment to reliable research support would make me a valuable addition to your team.

I would welcome the opportunity to contribute to your team with a combination of statistical training, programming experience in R, Python, C++, MATLAB, and Bash, and a strong commitment to reliable research support. Thank you for your time and consideration. I look forward to the opportunity to discuss how my background can support the department's research and analysis needs.

Sincerely,\\
\vspace{1em}

Nahid Hasan Sagar

\end{document}
